% P11 paper module: R14 exact polar-gauge separation and no-promotion countermodel.

\subsection{Exact polar-gauge separation and algebraic non-promotion}\label{subsec:o3m-polar-gauge}

Keep the notation of Remark~\ref{rem:polar-gauge} at fixed
$0<R<S<T_0<U$ on the odd sectors:
\[
 W:=W_{R,S,-}^{[T_0]},
 \qquad
 Q:=A_S^{1/2}WA_R^{-1/2},
\]
\[
 X_R=(G_{R,U}^-)^{1/2}(G_{R,T_0}^-)^{-1/2},
 \qquad
 X_S=(G_{S,U}^-)^{1/2}(G_{S,T_0}^-)^{-1/2}.
\]
The polar decompositions are
\[
 X_R=U_RA_R^{1/2},
 \qquad
 X_S=U_SA_S^{1/2}.
\]

\begin{proposition}[Uniqueness of the polar gauges]\label{prop:o3m-polar-unique}
The unitary factors $U_R,U_S$ are uniquely determined by the corresponding metric
pairs.  Explicitly,
\[
 \boxed{
 U_R=X_RA_R^{-1/2},
 \qquad
 U_S=X_SA_S^{-1/2}.
 }
\tag{O3M.1}\label{eq:o3m-polar-unique}
\]
Thus the gauge issue is not a nonuniqueness of polar decomposition.
\end{proposition}

\begin{proof}
For a boundedly invertible operator $X$, the positive factor in its polar decomposition
is uniquely $|X|=(X^*X)^{1/2}$ and the unitary factor is uniquely
$X|X|^{-1}$.  Here $|X_R|=A_R^{1/2}$ and $|X_S|=A_S^{1/2}$.
\end{proof}

Let
\[
 K_{R,S}^{T_0,U}:=W^*W_{R,S,-}^{[U]}
\]
be the cross-terminal kernel.  Recall the modulus defects
\[
 \mathscr K:=I-W^*Q,
 \qquad
 \mathscr N:=(I-WW^*)Q,
\]
so that
\[
 Q=W(I-\mathscr K)+\mathscr N.
\]

\begin{proposition}[Exact gauge/modulus separation of the cross kernel]
\label{prop:o3m-cross-separation}
One has
\[
 \boxed{
 K_{R,S}^{T_0,U}
 =\Bigl[W^*U_SW(I-\mathscr K)+W^*U_S\mathscr N\Bigr]U_R^*.
 }
\tag{O3M.2}\label{eq:o3m-cross-separation}
\]
In particular, with
\[
 \Gamma_U:=W^*U_SWU_R^*,
\]
\[
 \boxed{
 \|K_{R,S}^{T_0,U}-\Gamma_U\|
 \le \|\mathscr K\|+\|\mathscr N\|
 \le 2\|Q-W\|.
 }
\tag{O3M.3}\label{eq:o3m-gauge-error}
\]
Consequently, even under the hypothetical stronger assumption $Q-W\to0$, convergence
of the actual cross kernel to the identity still requires the separate gauge condition
\[
 \Gamma_U\to I.
\]
\end{proposition}

\begin{proof}
The exact polar-gauge identity~\eqref{eq:jd-polar-gauge} gives
\[
 K_{R,S}^{T_0,U}=W^*U_SQU_R^*.
\]
Insert $Q=W(I-\mathscr K)+\mathscr N$ to obtain
\eqref{eq:o3m-cross-separation}.  Since $W$ is an isometry and $U_S,U_R$ are unitary,
\[
 \|K_{R,S}^{T_0,U}-\Gamma_U\|
 \le\|\mathscr K\|+\|\mathscr N\|.
\]
For every source vector $f$, the two summands in
$Q-W=-W\mathscr K+\mathscr N$ are orthogonal.  Hence
$\|\mathscr K\|\le\|Q-W\|$ and
$\|\mathscr N\|\le\|Q-W\|$, proving~\eqref{eq:o3m-gauge-error}.
\end{proof}

The next finite-dimensional model shows that no conclusion about the actual transport
can be extracted from the modulus data alone, even in the extreme case
$Q=W$ and $\Theta=0$ exactly.

\begin{proposition}[Exact algebraic countermodel to modulus-only promotion]
\label{prop:o3m-countermodel}
There exist finite-dimensional positive metric data satisfying the exact pullback,
relative-compression and polar-gauge identities for which
\[
 \boxed{Q=W=I,\qquad \Theta=0,}
\tag{O3M.4}\label{eq:o3m-perfect-modulus}
\]
but
\[
 \boxed{W^{[U]}\ne W,\qquad K^{T_0,U}\ne I.}
\tag{O3M.5}\label{eq:o3m-nontrivial-gauge}
\]
Moreover, with one fixed baseline system one can choose a sequence of future positive
metrics for which~\eqref{eq:o3m-perfect-modulus} holds at every step while the actual
future transports are not Cauchy.
\end{proposition}

\begin{proof}
Work on $\mathbb C^2$ and let
\[
 B=\begin{pmatrix}2&1\\1&2\end{pmatrix}>0,
 \qquad
 A=\begin{pmatrix}4&0\\0&1\end{pmatrix}>0.
\]
These matrices do not commute.  Put
\[
 G_{R,T_0}=I,
 \qquad
 G_{S,T_0}=B,
 \qquad
 J=B^{-1/2},
\]
and at the future horizon put
\[
 G_{R,U}=A,
 \qquad
 G_{S,U}=B^{1/2}AB^{1/2}.
\]
Then
\[
 J^*G_{S,T_0}J=G_{R,T_0},
 \qquad
 J^*G_{S,U}J=G_{R,U},
\]
so the exact terminal pullback identities hold.  The normalized base transport is
\[
 W=G_{S,T_0}^{1/2}JG_{R,T_0}^{-1/2}=I.
\]
The relative metrics are
\[
 A_R=G_{R,T_0}^{-1/2}G_{R,U}G_{R,T_0}^{-1/2}=A,
\]
\[
 A_S=G_{S,T_0}^{-1/2}G_{S,U}G_{S,T_0}^{-1/2}=A.
\]
Therefore
\[
 Q=A_S^{1/2}WA_R^{-1/2}=I,
\]
and the Jensen defect vanishes identically: $\Theta=0$.

On the other hand the actual future normalized transport is
\[
 V:=W^{[U]}
 =(B^{1/2}AB^{1/2})^{1/2}B^{-1/2}A^{-1/2}.
\tag{O3M.6}\label{eq:o3m-V}
\]
It is unitary by the pullback identity.  If $V=I$, then
\[
 (B^{1/2}AB^{1/2})^{1/2}=A^{1/2}B^{1/2}.
\]
The left side is positive selfadjoint, so the right side would be selfadjoint.  Hence
$A^{1/2}$ and $B^{1/2}$ would commute, equivalently $A$ and $B$ would commute, a
contradiction.  Thus $V\ne I$, proving~\eqref{eq:o3m-nontrivial-gauge}.

For the sequential statement keep the same fixed baseline data $B,J$ and alternate the
future source metric between $A$ and $I$.  Define the target future metric each time by
$B^{1/2}A_nB^{1/2}$.  At every step one still has $A_R=A_S=A_n$, hence
$Q=W=I$ and $\Theta=0$.  For $A_n=I$ the actual future transport is $I$, whereas for
$A_n=A$ it is the fixed nonidentity unitary $V$ in~\eqref{eq:o3m-V}.  The resulting
future transports therefore alternate between two distinct operators and are not
Cauchy.
\end{proof}

\begin{remark}[R14 firewall]\label{rem:o3m-firewall}
Proposition~\ref{prop:o3m-countermodel} is an algebraic non-promotion theorem: the exact
P11 identities by themselves do not imply terminal convergence from modulus/Jensen
data, even under the idealized condition $\Theta=0$.  It is not a counterexample to the
actual P11 metric family.  The actual asymptotic behavior of $U_R,U_S$ and
$K_{R,S}^{T_0,U}$ remains open and may still be constrained by additional arithmetic or
source structure not present in the abstract model.

In particular, the divergence
$\chi\|\Theta\|\to\infty$ proved in
Theorem~\ref{thm:o3l-polynomial-witness} only rules out the auxiliary sufficient
Jensen-product route.  It neither proves that $Q-W$ fails to tend to zero nor that the
actual cross-terminal kernel fails to tend to the identity.
\end{remark}
